\documentclass[11pt,a4paper]{article}
\usepackage[utf8]{inputenc}
\usepackage[T1]{fontenc}
\usepackage[margin=1in]{geometry}
\usepackage{xcolor}
\usepackage{hyperref}

\definecolor{primarycolor}{RGB}{33,37,41}
\definecolor{secondarycolor}{RGB}{0,102,204}
\hypersetup{colorlinks=true,linkcolor=secondarycolor,urlcolor=secondarycolor}
\pagestyle{empty}
\setlength{\parskip}{0.9em}
\setlength{\parindent}{0pt}

\begin{document}

Dear Verdane Care Hiring Team,

I am applying for the Full-Stack Engineer role because it brings together patient-facing web development and a strong commitment to accessibility. Verdane Care’s focus on making care feel simple, along with its WCAG 2.2 AA commitment, particularly appeals to me.

As a third-year Computing student at Liffey Institute of Technology, I have developed accessible React and TypeScript interfaces through a software development internship and several projects. At Harbourline Media, I implemented accessibility fixes and React components for a news site, including a keyboard- and screen-reader-friendly image gallery documented in Storybook. In my Tidewatch final-year project, I tested the interface with NVDA and VoiceOver and achieved a Lighthouse accessibility score of 100.

Tidewatch also gave me project experience across the stack: I built a Node.js/Express API that caches data from public weather and tide APIs, and used PostgreSQL for user favourites. I have also used Jest, and my freelance work building a yoga studio booking site gave me experience with React and Next.js. As a teaching assistant in a first-year web development lab, I helped approximately 30 students over one semester.

I would welcome the opportunity to discuss how my accessibility work, frontend experience, and project-level backend development could contribute to Verdane Care’s patient experience features. Thank you for considering my application.

Yours sincerely,\\
Aiyana Brooks-Tan

\end{document}